% !TeX root = ../../../main.tex
% chapters/sections/chapter5/subsections/two_plane/classification.tex

\subsection{方程式体系に含まれる式の分類}
\label{sec:chapter5_two_plane_classification}
本節では、代表的家計の方程式体系に含まれる全式を「供給群」「需要群」「外部群」の3つの群に分類する。
後述のように、
主に供給群の式を統合することによりYS（生産供給）曲線およびCS（消費供給）曲線が導出され、
主に需要群の式を統合することによりYD（生産需要）曲線およびCD（消費需要）曲線が導出される。
YS曲線とCS曲線は外部群の式を統合することにより互いに他に変換可能であり、
YD曲線とCD曲線も外部群の式を統合することにより互いに他に変換可能である。
このようにして導かれたYS曲線とYD曲線を用いて\((y_t^H,p_t^H)\)平面における均衡分析がおこなわれる。
同様にCS曲線とCD曲線を用いて\((c_t^{H\to W},p_t^{H\to W})\)平面における均衡分析がおこなわれる。

\subsubsection*{1.供給群}
\label{sec:chapter5_two_plane_classification_supply}
財市場における供給側の均衡条件を表す方程式群であり、主にこれらを統合することによりYS曲線が導かれる。
\begin{itemize}
    \item
    \label{itm:chapter5_two_plane_classification_supply_p_H_tilde}
        \textbf{最適価格\(\widetilde{p}_t^H\)の定義式}:
        式
        \eqref{form:chapter3_representative_household_phillips_home_representative_p_H_tilde_def}
        により定義された。
        \begin{equation}
        \label{form:chapter5_two_plane_classification_supply_p_H_tilde_def}
            (\widetilde{p}_t^H)^{1+\theta^H}=\frac{\theta^H}{\theta^H-1}\frac{v_t^H}{w_t^H}
        \end{equation}
    \item
    \label{itm:chapter5_two_plane_classification_supply_v_H}
        \textbf{最適価格\(\widetilde{p}_t^H\)の補助変数\(v_t^H\)の式}:
        式
        \eqref{form:chapter3_representative_household_phillips_home_representative_v_H_eq}
        により導かれた。
        \begin{equation}
        \label{form:chapter5_two_plane_classification_supply_v_H_eq}
            v_t^H=\phi^H\frac{(y_t^H N)^2}{(a_t^H)^2}(p_t^H)^{2\theta^H}+\beta_t^H\xi^H\operatorname{E}_t[v_{t+1}^H]
        \end{equation}
    \item
    \label{itm:chapter5_two_plane_classification_supply_w_H}
        \textbf{最適価格\(\widetilde{p}_t^H\)の補助変数\(w_t^H\)の式}:
        式
        \eqref{form:chapter3_representative_household_phillips_home_representative_w_H_eq}
        により導かれた。
        \begin{equation}
        \label{form:chapter5_two_plane_classification_supply_w_H_eq}
            w_t^H=\lambda_t^H(1-\tau_t^H)y_t^H N(p_t^H)^{\theta^H}+\beta_t^H\xi^H\operatorname{E}_t[w_{t+1}^H]
        \end{equation}
    \item
    \label{itm:chapter5_two_plane_classification_supply_p_H_bar}
        \textbf{正規化された生産者物価指数(PPI)\(\bar{p}_t^H\)の式}:
        式
        \eqref{form:chapter3_representative_household_phillips_home_representative_p_H_bar_dynamics_p_H_bar_modified2_eq}
        により示された。
        \begin{equation}
        \label{form:chapter5_two_plane_classification_supply_p_H_bar_eq}
            (\bar{p}_t^H)^{1-\theta^H}=(1-\xi^H)(\widetilde{p}_t^H)^{1-\theta^H}+\xi^H(\bar{p}_{t-1}^H)^{1-\theta^H}
        \end{equation}
    \item
    \label{itm:chapter5_two_plane_classification_supply_p_H_bar_p_H}
        \textbf{正規化された生産者物価指数(PPI)\(\bar{p}_t^H\)と生産者物価指数(PPI)\(p_t^H\)の関係式}:
        式
        \eqref{form:chapter3_households_stage1a_normalized_ppi_definitions_p_H_bar_eq}
        により定義された。
        \begin{equation}
        \label{form:chapter5_two_plane_classification_supply_p_H_bar_p_H_eq}
            \bar{p}_t^H=N^{-\frac{1}{1-\theta^H}}p_t^H
        \end{equation}
    \item
    \label{itm:chapter5_two_plane_classification_supply_Delta_H}
        \textbf{価格分散\(\Delta_t^H\)の式}:
        式
        \eqref{form:chapter3_representative_household_phillips_home_representative_Delta_H_dynamics_Delta_H_modified2_eq}
        により示された。
        \begin{equation}
        \label{form:chapter5_two_plane_classification_supply_Delta_H_eq}
            \Delta_t^H=(1-\xi^H)N\left(\frac{\widetilde{p}_t^H}{p_t^H}\right)^{-\theta^H}+\xi^H\left(\frac{p_t^H}{p_{t-1}^H}\right)^{\theta^H}\Delta_{t-1}^H
        \end{equation}
    \item
    \label{itm:chapter5_two_plane_classification_supply_y_H}
        \textbf{生産関数}:
        式
        \eqref{form:chapter3_representative_household_production_home_original_y_H_modified_eq}
        より示された。
        \begin{equation}
        \label{form:chapter5_two_plane_classification_supply_y_H_eq}
            y_t^H=\frac{a_t^H l_t^H}{\Delta_t^H}
        \end{equation}
    \item
    \label{itm:chapter5_two_plane_classification_external_pi_H}
        \textbf{生産者物価指数\(p_t^H\)のグロス・インフレ率\(\pi_t^H\)の定義式}:
        式
        \eqref{form:chapter3_policies_monetary_pi_H_def}
        により定義された。
        \begin{equation}
        \label{form:chapter5_two_plane_classification_external_pi_H_def}
            \pi_t^H\equiv\frac{p_t^H}{p_{t-1}^H}
        \end{equation}
\end{itemize}

\subsubsection*{2.需要群}
\label{sec:chapter5_two_plane_classification_demand}
財市場における需要側の均衡条件を表す方程式群であり、主にこれら統合することによりCD曲線が導かれる。
\begin{itemize}
    \item
    \label{itm:chapter5_two_plane_classification_demand_lambda_H_i_H}
        \textbf{自国コンティンジェント債券に関するオイラー方程式}:
        式
        \eqref{form:chapter3_representative_household_euler_contingent_home_representative_lambda_H_i_H_eq}
        より示された。
        \begin{equation}
        \label{form:chapter5_two_plane_classification_demand_lambda_H_i_H_eq}
            \lambda_t^H=\beta_t^H(1+i_t^H)\operatorname{E}_t[\lambda_{t+1}^H]
        \end{equation}   
    \item
    \label{itm:chapter5_two_plane_classification_demand_lambda_H_p_H_W_c_H_W}
        \textbf{所得の限界効用\(\lambda_t^H\)の式}:
        式
        \eqref{form:chapter3_representative_household_lambda_home_representative_lambda_H_p_H_W_c_H_W_eq}
        により示された。
        \begin{equation}
        \label{form:chapter5_two_plane_classification_demand_lambda_H_p_H_W_c_H_W_eq}
            \lambda_t^H=\frac{1}{p_t^{H\to W}c_t^{H\to W}}
        \end{equation}
\end{itemize}

\subsubsection*{3.外部群}
\label{sec:chapter5_two_plane_classification_connection}
財市場以外における均衡条件を表す方程式群である。
これらを合わせることによりYS曲線とCS曲線を互いに他に変換でき、YD曲線とCD曲線も互いに他に変換できる。
\begin{itemize}
    \item
    \label{itm:chapter5_two_plane_classification_external_euler}
        \textbf{外国通貨建てリスクフリー債券のオイラー方程式}:
        式
        \eqref{form:chapter3_representative_household_euler_risk_free_home_representative_lambda_H_i_F_eq}
        により示された。
        \begin{equation}
        \label{form:chapter5_two_plane_classification_external_lambda_H_i_F_eq}
            \lambda_t^H e_t^{/*}=(1+i_t^F)\beta_t^H\operatorname{E}_t[\lambda_{t+1}^H e_{t+1}^{/*}]
        \end{equation}
    \item
    \label{itm:chapter5_two_plane_classification_external_lambda}
        \textbf{所得の限界効用\(\lambda_t^H\)の式}:
        式
        \eqref{form:chapter3_representative_household_lambda_home_representative_lambda_H_p_H_W_c_H_W_eq}
        により示された。
        \begin{equation}
        \label{form:chapter5_two_plane_classification_external_lambda_H_p_H_W_c_H_W_eq}
            \lambda_t^H=\frac{1}{p_t^{H\to W}c_t^{H\to W}}
        \end{equation}
    \item
    \label{itm:chapter5_two_plane_classification_external_lambda_H_star_def}
        \textbf{外国通貨所得の限界効用\(\lambda_t^{H/*}\)の定義式}:
        付録
        \ref{chap:appendix_insulation_effect}
        において定義された、名目為替レートと所得の限界効用を結合する定義式である。
        \begin{equation}
        \label{form:chapter5_two_plane_classification_external_lambda_H_star_def_eq}
            \lambda_t^{H/*}\equiv e_t^{/*} \lambda_t^H
        \end{equation}
    \item
    \label{itm:chapter5_two_plane_classification_external_demand_home_goods}
        \textbf{自国財消費指数\(c_t^{H\to H}\)の需要関数}:
        式
        \eqref{form:chapter3_representative_household_demand_home_representative_c_H_H_eq}
        により示された。
        \begin{equation}
        \label{form:chapter5_two_plane_classification_external_c_H_H_eq}
            c_t^{H\to H}=\alpha^H\frac{p_t^{H\to W}}{p_t^H}c_t^{H\to W}
        \end{equation}
    \item
    \label{itm:chapter5_two_plane_classification_external_demand_foreign_goods}
        \textbf{外国財消費指数\(c_t^{H\to F}\)の需要関数}:
        式
        \eqref{form:chapter3_representative_household_demand_home_representative_c_H_F_eq}
        により示された。
        \begin{equation}
        \label{form:chapter5_two_plane_classification_external_c_H_F_eq}
            c_t^{H\to F}=(1-\alpha^H)\frac{p_t^{H\to W}}{p_t^F}c_t^{H\to W}
        \end{equation}
    \item
    \label{itm:chapter5_two_plane_classification_external_lop}
        \textbf{一物一価の法則(LOP)}:
        式
        \eqref{form:chapter3_representative_household_lop_home_representative_p_H_eq}
        により仮定された。
        \begin{equation}
        \label{form:chapter5_two_plane_classification_external_p_H_eq}
            p_t^H=e_t^{/*}p_t^{H*}
        \end{equation}
    \item
    \label{itm:chapter5_two_plane_classification_external_cpi}
        \textbf{消費者物価指数(CPI)\(p_t^{H\to W}\)の式}:
        式
        \eqref{form:chapter3_households_stage1b_cpi_descriptions_p_H_W_eq}
        により示された。
        \begin{equation}
        \label{form:chapter5_two_plane_classification_external_p_H_W_eq}
            p_t^{H\to W}\equiv(p_t^H)^{\alpha^H}(p_t^F)^{1-\alpha^H}
        \end{equation}
    \item
    \label{itm:chapter5_two_plane_classification_external_constraint}
        \textbf{資源制約式}:
        式
        \eqref{form:chapter3_representative_household_resource_constraints_home_representative_constraint_eq}
        により示された。
        \begin{equation}
        \label{form:chapter5_two_plane_classification_external_constraint_eq}
            b_{t+1}^H+p_t^{H\to W}c_t^{H\to W}=p_t^H y_t^H+(1+i_{t-1}^F)\frac{e_t^{/*}}{e_{t-1}^{/*}}b_t^H
        \end{equation}
    \item
    \label{itm:chapter5_two_plane_classification_external_goods_market}
        \textbf{財市場の均衡条件}:
        式
        \eqref{form:chapter3_representative_household_goods_market_home_representative_y_H_eq}
        により示された。
        \begin{equation}
        \label{form:chapter5_two_plane_classification_external_y_H_eq}
            y_t^H=c_t^{H\to H}+\frac{M}{N}c_t^{F\to H}
        \end{equation}
    \item
    \label{itm:chapter5_two_plane_classification_external_international_bond_equilibrium}
        \textbf{国際リスクフリー債券市場の均衡条件}:
        式
        \eqref{form:chapter3_assumptions_bond_markets_international_equilibrium}
        により示された、世界全体の債券市場における均衡条件である。
        \begin{equation}
        \label{form:chapter5_two_plane_classification_external_international_bond_equilibrium_eq}
            \sum_{h\in H}b_t^{h\to F}+\sum_{f\in F}e_t^{/*}b_t^{f\to H*}=0
        \end{equation}
    \item
    \label{itm:chapter5_two_plane_classification_external_steady_state}
        \textbf{初期定常状態における外国通貨建てリスクフリー債券\(b_s^H\)の無保有}:
        式
        \eqref{form:chapter3_representative_household_steady_state_home_representative_b_s_H_eq}
        により示された。
        \begin{equation}
        \label{form:chapter5_two_plane_classification_external_b_s_H_eq}
            b_s^H=0
        \end{equation}
    \item
    \label{itm:chapter5_two_plane_classification_external_foreign_system}
        \textbf{外国の方程式系に含まれる全方程式}
\end{itemize}